% [1.5 pts] ~0.6 páginas. Recomendación concreta y accionable; no repetir el análisis.
\section{Conclusiones y recomendaciones}
\begin{itemize}
  \item Se recomienda cubrir con derivados los dos riesgos de mercado que hoy se asumen sin cobertura y que, con
  utilidad bruta casi nula y liquidez restringida, comprometen la continuidad operativa: la posición pasiva neta de
  S/~\PenNetoAbs\ millones, sin amortizaciones del BN hasta 2027, y el inventario de crudo, expuesto a un WTI que cayó
  \WTICaidaVeinticinco\ en 2025 y hoy cotiza en US\$~\WTISpot/bl con una volatilidad implícita de \VolCrudo.
  \item \textbf{Estrategia 1:} el CCS amortizable sobre el préstamo del BN, pagando en dólares al \CCSTasaAllIn\ (con
  CVA), junto con NDF de compra de soles, cubre el \CoberturaFX\ de la exposición en soles y reduce la pérdida ante
  una caída de 10\,\% del tipo de cambio de US\$~\EscSinCobMenosDiez\ millones a US\$~\EscConCobMenosDiez\ millones, sin
  costo inicial y con contabilidad de coberturas para el CCS; su costo es el diferencial de tasas (\NDFCostoAnual\ anual).
  \item \textbf{Estrategia 2:} cubrir el \CoberturaCrudoPct\ del inventario de crudo, por partes iguales con un swap de WTI a
  US\$~\SwapPrecio/bl y con collares de costo cero en tres capas mensuales (puts al \PutPctFwd\ de cada futuro), reduce la
  pérdida ante una caída de 30\,\% del WTI a US\$~\PerdidaMaxMixto\ millones (frente a US\$~\PerdidaSinCobTreinta\
  millones) sin pagar prima.
  \item La restricción de liquidez obliga a preferir instrumentos OTC sin costo inicial en lugar de futuros con
  márgenes diarios; aun así, en un estrés conjunto el colateral llegaría a US\$~\ColateralConjunto\ millones, por lo que
  la ejecución se condiciona a un umbral CSA de al menos US\$~\UmbralNecesario\ millones o a la garantía del Estado. Para el riesgo de tasa conviene esperar al diseño de la refinanciación
  y evaluar entonces un swap de inicio diferido, solo si la refinanciación es altamente probable.
  \item Se recomienda a la Gerencia General: (1) presentar la Política de Cobertura al Directorio; (2) iniciar la
  negociación ISDA/CSA con al menos tres bancos; (3) obtener la opinión legal y gestionar el respaldo del Estado (MINEM, D.U.~\DUCrisis) para las líneas de
  derivados; y (4) ejecutar primero el CCS, que cubre el riesgo de mayor impacto, una vez firmado el CSA con el umbral
  requerido.
\end{itemize}
